% ЗГЕНЕРОВАНО scripts/build_format_v2_chapter.py — не редагувати вручну.
% Джерело: editorial/cards/19.*.json, format-v2/rozdil-19.json, paraleli.json.

\naukchapter{19}{8}{Глава 19. Де це все?}
\naukintro{Девʼятнадцята глава складається з восьми віршів, і сім із них мають одну форму: питання де те й де те — і відповідь у мене, що стоїть у власній величі. Це не заперечення перелічених речей: текст питає, де вони, у того, хто вже стоїть. Восьмий вірш обриває навіть цю форму: досить розмов.}
\vsec{19.1}
\vdeva{तत्त्वविज्ञानसन्दंशम् आदाय हृदयोदरात् । नानाविधपरामर्शशल्योद्धारः कृतो मया ॥}
\vtranslit{таттвавіджнянасандамшам адая хрідайодарат | нанавідгапарамаршашалйоддгарах кріто мая ||}
\vperek{Взявши кліщі знання істини, я витяг із нутра серця колючки різноманітних хибних уявлень.}
\vcom{Построкове зіставлення: редакторська чернетка, не рецензовано.}
\vcom{Рядок а}
\vdeva{तत्त्वविज्ञानसन्दंशम् आदाय हृदयोदरात्}
\vtranslit{таттвавіджнянасандамшам адая хрідайодарат}
\vperek{Взявши кліщі знання істини,}
\vcom{tattvavijñāna («знання істини») входить до складного слова із sandaṃśa («щипці, кліщі»); ādāya — дієприслівникове «взявши». hṛdayodarāt означає «з нутра серця», але в українському цілому реченні цей зворот стоїть у другій частині: поділ перекладу є смисловим, не послівним.}
\vcom{Рядок б}
\vdeva{नानाविधपरामर्शशल्योद्धारः कृतो मया}
\vtranslit{нанавідгапарамаршашалйоддгарах кріто мая}
\vperek{я витяг із нутра серця колючки різноманітних хибних уявлень.}
\vcom{nānāvidha — «різноманітний», parāmarśa тут стосується схоплювання або уявлень, śalya — «колючка». uddhāraḥ kṛto mayā передає виконане «мною» вилучення; український активний вираз «я витяг» зберігає чинну редакцію, але не відтворює буквально пасивну побудову санскриту.}
\vcom{Вірш — не декларація, а опис процедури, і метафора тут хірургічна: sandaṃśa — кліщі, якими виймають стороннє тіло з рани. Три слова працюють разом: śalya («колючка») — те, що застрягло; uddhāra («витягання») — дія; hṛdayodara («нутро серця») — місце, і воно назване нутром, а не поверхнею. parāmarśa тут не «думка», а хибне схоплювання — розумове прийняття чогось за інше. Тобто йдеться не про прибирання думок, а про виймання того, що застрягло й болить. РЕФРЕН ГЛАВИ. 19.2–19.7 закінчуються однаково: svamahimni sthitasya me («у мене, що стоїть у власній величі»). Кожен із цих шести віршів — низка питань kva («де»), і жодне не дістає відповіді. 19.1 і 19.8 стоять поза серією: перший описує, ЯК це сталося, останній — ЩО з цього випливає для мови.}
\vpar{Українська паралель (наша інтерпретація)}{Григорій Сковорода, «Асхань. Симфонія, наречена книга Асхань, про пізнання самого себе».}{Наша паралель: у Сковороди пізнання себе теж постає як виймання зайвого, а не як набирання нового. Наше порівняння — типологічне, не твердження про джерело.}
\vfig{ilustratsii/19-1-thread.pdf}{Кліщі знання витягають колючки з серця}{Схема наша, не історична ілюстрація.}

\vsec{19.2}
\vdeva{क्व धर्मः क्व च वा कामः क्व चार्थः क्व विवेकता । क्व द्वैतं क्व च वाद्वैतं स्वमहिम्नि स्थितस्य मे ॥}
\vtranslit{ква дгармах ква ча ва камах ква чартхах ква вівеката | ква двайтам ква ча вадвайтам свамахімні стхітася ме ||}
\vperek{Де дгарма, де бажання, де користь, де вміння розрізняти? Де двоїстість, де недвоїстість — у мене, що стоїть у власній величі?}
\vcom{Построкове зіставлення: редакторська чернетка, не рецензовано.}
\vcom{Рядок а}
\vdeva{क्व धर्मः क्व च वा कामः क्व चार्थः क्व विवेकता}
\vtranslit{ква дгармах ква ча ва камах ква чартхах ква вівеката}
\vperek{Де дгарма, де бажання, де користь, де вміння розрізняти?}
\vcom{Повтор kva («де?») об'єднує dharmaḥ, kāmaḥ, arthaḥ і термін розрізнення. Чинна картка має vivekatā, тоді як електронний GRETIL читає vivekitā: цю лексичну відмінність потрібно окремо перевірити; механічно виправляти авторський український переклад не можна.}
\vcom{Рядок б}
\vdeva{क्व द्वैतं क्व च वाद्वैतं स्वमहिम्नि स्थितस्य मे}
\vtranslit{ква двайтам ква ча вадвайтам свамахімні стхітася ме}
\vperek{Де двоїстість, де недвоїстість — у мене, що стоїть у власній величі?}
\vcom{Два питання kva dvaitaṃ і kva ... advaitaṃ охоплюють і двоїстість, і недвоїстість. svamahimni — місцевий відмінок «у власній величі», sthitasya me — родове «мене, що стоїть»; формула повторюється у віршах 19.2–19.7, і її український відповідник збережено.}
\vcom{Перший із шести віршів-переліків. Три з перших чотирьох понять — це три з чотирьох пурушартг (дгарма, користь, бажання); четверте, визволення, не назване, бо його місце займає сама позиція того, хто питає. Найважливіше — остання пара: питається не лише «де двоїстість», а й «де недвоїстість». Тобто заперечується не світ на користь єдності, а сама вісь, на якій обидва поняття мають сенс. svamahimni («у власній величі») — не гордість, а міра: те, що не має ні межі, ні місця. РЕФРЕН ГЛАВИ. 19.2–19.7 закінчуються однаково: svamahimni sthitasya me («у мене, що стоїть у власній величі»). Кожен із цих шести віршів — низка питань kva («де»), і жодне не дістає відповіді. 19.1 і 19.8 стоять поза серією: перший описує, ЯК це сталося, останній — ЩО з цього випливає для мови.}
\vpar{Українська паралель (наша інтерпретація)}{Митрополит Іларіон Київський, «Слово про закон і благодать» (XI ст.).}{Наша паралель: у київському Слові перелік того, що вже не важить, теж будується як питання, а не як твердження. Наше порівняння — типологічне; датування цієї книги науково не встановлене.}
\vfig{ilustratsii/19-2-neti.pdf}{Дгарма, бажання й двоїстість відпали}{Схема наша, не історична ілюстрація.}

\vsec{19.3}
\vdeva{क्व भूतं क्व भविष्यद् वा वर्तमानम् अपि क्व वा । क्व देशः क्व च वा नित्यं स्वमहिम्नि स्थितस्य मे ॥}
\vtranslit{ква бгутам ква бгавішяд ва вартаманам апі ква ва | ква дешах ква ча ва нітям свамахімні стхітася ме ||}
\vperek{Де минуле, де майбутнє, де навіть теперішнє? Де сторона світу, де вічність — у мене, що стоїть у власній величі?}
\vcom{Построкове зіставлення: редакторська чернетка, не рецензовано.}
\vcom{Рядок а}
\vdeva{क्व भूतं क्व भविष्यद् वा वर्तमानम् अपि क्व वा}
\vtranslit{ква бгутам ква бгавішяд ва вартаманам апі ква ва}
\vperek{Де минуле, де майбутнє, де навіть теперішнє?}
\vcom{bhūtaṃ, bhaviṣyad і vartamānam позначають минуле, майбутнє та теперішнє; api посилює згадку про теперішнє («навіть»). Питальне kva повторюється, а не називає фактичної відсутності календарного часу в житті читача.}
\vcom{Рядок б}
\vdeva{क्व देशः क्व च वा नित्यं स्वमहिम्नि स्थितस्य मे}
\vtranslit{ква дешах ква ча ва нітям свамахімні стхітася ме}
\vperek{Де сторона світу, де вічність — у мене, що стоїть у власній величі?}
\vcom{deśaḥ — місце або сторона, nityaṃ — постійне чи вічне. Рефрен svamahimni sthitasya me співвідносить простір і тривалість зі станом мовця; це філософське питання, не фізичне заперечення простору.}
\vcom{Три часи названі поспіль, і слово api («навіть») стоїть саме при теперішньому: минуле й майбутнє відсутні легко, а теперішнє — те, що здається найбільш незаперечним. Друга пара — deśa («місце») і nitya («вічне») — зіставляє простір із часом; разом обидві пари охоплюють усе, у чому можна було б розміститися. Вірш показує, що «стояти у власній величі» означає не займати іншого місця, а не мати місця взагалі. РЕФРЕН ГЛАВИ. 19.2–19.7 закінчуються однаково: svamahimni sthitasya me («у мене, що стоїть у власній величі»). Кожен із цих шести віршів — низка питань kva («де»), і жодне не дістає відповіді. 19.1 і 19.8 стоять поза серією: перший описує, ЯК це сталося, останній — ЩО з цього випливає для мови.}
\vpar{Українська паралель (ПОТРЕБУЄ ПЕРЕВІРКИ)}{українська філософська лірика, не визначено.}{ПОТРЕБУЄ ПЕРЕВІРКИ: питання, де минуле й майбутнє, близьке до української філософської лірики, але конкретного твору ми не звірили.}
\vfig{ilustratsii/19-3-neti.pdf}{Минуле, майбутнє й вічність відпали}{Схема наша, не історична ілюстрація.}

\vsec{19.4}
\vdeva{क्व चात्मा क्व च वानात्मा क्व शुभं क्वाशुभं तथा । क्व चिन्ता क्व च वाचिन्ता स्वमहिम्नि स्थितस्य मे ॥}
\vtranslit{ква чатма ква ча ванатма ква шубгам квашубгам татха | ква чінта ква ча вачінта свамахімні стхітася ме ||}
\vperek{Де Атман, де не-Атман, де добре, де недобре? Де думка, де відсутність думки — у мене, що стоїть у власній величі?}
\vcom{Построкове зіставлення: редакторська чернетка, не рецензовано.}
\vcom{Рядок а}
\vdeva{क्व चात्मा क्व च वानात्मा क्व शुभं क्वाशुभं तथा}
\vtranslit{ква чатма ква ча ванатма ква шубгам квашубгам татха}
\vperek{Де Атман, де не-Атман, де добре, де недобре?}
\vcom{Пари ātmā / anātmā та śubhaṃ / aśubhaṃ запитально зіставлені через kva. Чинна картка бере SARIT kvāśubhaṃ tathā; GRETIL має kvāsubhaṃ yathā — відмінність торкається і слова, і зв'язки. Жодного з цих електронних читань не оголошено незалежно підтвердженим рукописним оригіналом.}
\vcom{Рядок б}
\vdeva{क्व चिन्ता क्व च वाचिन्ता स्वमहिम्नि स्थितस्य मे}
\vtranslit{ква чінта ква ча вачінта свамахімні стхітася ме}
\vperek{Де думка, де відсутність думки — у мене, що стоїть у власній величі?}
\vcom{cintā означає думку або турботу, acintā — її відсутність; запитання поширюється на обидва боки пари. Рефрен svamahimni sthitasya me збережено однаково з сусідніми віршами. Цей рядок не дає підстав наказувати українському читачеві відмовитися від мислення.}
\vcom{Три пари, і в кожній заперечується не один бік, а обидва. Найгостріша — перша: питати «де Атман» і «де не-Атман» однаково означає не знайти місця ні для того, ні для того, тобто сама пара втрачає ґрунт. Друга пара (śubha / aśubha) — та сама, що в 13.7, де вона відкидалася цілою. Третя (cintā / acintā) важлива тим, що сюди потрапляє й «відсутність думки»: навіть спокій розуму не стає позицією, з якої можна стояти. РЕФРЕН ГЛАВИ. 19.2–19.7 закінчуються однаково: svamahimni sthitasya me («у мене, що стоїть у власній величі»). Кожен із цих шести віршів — низка питань kva («де»), і жодне не дістає відповіді. 19.1 і 19.8 стоять поза серією: перший описує, ЯК це сталося, останній — ЩО з цього випливає для мови.}
\vpar{Українська паралель (наша інтерпретація)}{Олександр Потебня, «Думка й мова» (1862).}{Наша паралель: Потебня теж показує, що пара думка й відсутність думки є справою самого мірила, а не речей. Наше порівняння — типологічне.}
\vfig{ilustratsii/19-4-neti.pdf}{Атман і не-Атман, думка й не-думка відпали}{Схема наша, не історична ілюстрація.}

\vsec{19.5}
\vdeva{क्व स्वप्नः क्व सुषुप्तिर् वा क्व च जागरणं तथा । क्व तुरीयं भयं वापि स्वमहिम्नि स्थितस्य मे ॥}
\vtranslit{ква свапнах ква сушуптір ва ква ча джаґаранам татха | ква туріям бгаям вапі свамахімні стхітася ме ||}
\vperek{Де сон зі сновидіннями, де глибокий сон, де неспання? Де четвертий стан, де навіть страх — у мене, що стоїть у власній величі?}
\vcom{Построкове зіставлення: редакторська чернетка, не рецензовано.}
\vcom{Рядок а}
\vdeva{क्व स्वप्नः क्व सुषुप्तिर् वा क्व च जागरणं तथा}
\vtranslit{ква свапнах ква сушуптір ва ква ча джаґаранам татха}
\vperek{Де сон зі сновидіннями, де глибокий сон, де неспання?}
\vcom{svapnaḥ — сон зі сновидіннями, suṣuptiḥ — глибокий сон, jāgaraṇaṃ — неспання. kva ставить під питання всі три стани; така трійця має філософський контекст, але тут не є описом медичного або психологічного дослідження сну.}
\vcom{Рядок б}
\vdeva{क्व तुरीयं भयं वापि स्वमहिम्नि स्थितस्य मे}
\vtranslit{ква туріям бгаям вапі свамахімні стхітася ме}
\vperek{Де четвертий стан, де навіть страх — у мене, що стоїть у власній величі?}
\vcom{turīyaṃ — «четвертий» стан, bhayaṃ — «страх», api посилює друге питання. Відмінність turīyaṃ / turiyaṃ між електронними текстами потребує орфографічного обліку; висновок про рукописне різночитання з неї не випливає.}
\vcom{Три стани — неспання, сон зі сновидіннями, глибокий сон — це стандартна схема упанішад, і четвертий (turīya) стоїть поза нею. Тут відкидається й він: навіть те, що мало б бути виходом, не має місця. Страх (bhaya) доданий останнім, і це не випадково: страх — те, що тримає в межах трьох станів, і якщо його немає, то й виходити нікуди. РЕФРЕН ГЛАВИ. 19.2–19.7 закінчуються однаково: svamahimni sthitasya me («у мене, що стоїть у власній величі»). Кожен із цих шести віршів — низка питань kva («де»), і жодне не дістає відповіді. 19.1 і 19.8 стоять поза серією: перший описує, ЯК це сталося, останній — ЩО з цього випливає для мови.}
\vpar{Українська паралель (ПОТРЕБУЄ ПЕРЕВІРКИ)}{українська релігійна проза, не визначено.}{ПОТРЕБУЄ ПЕРЕВІРКИ: перелік станів від сну до четвертого стану найближчий до української релігійної прози, але назвати твір без звірки ми не можемо.}
\vfig{ilustratsii/19-5-neti.pdf}{Сон, неспання й четвертий стан відпали}{Схема наша, не історична ілюстрація.}

\vsec{19.6}
\vdeva{क्व दूरं क्व समीपं वा बाह्यं क्वाभ्यन्तरं क्व वा । क्व स्थूलं क्व च वा सूक्ष्मं स्वमहिम्नि स्थितस्य मे ॥}
\vtranslit{ква дурам ква саміпам ва бахям квабгянтарам ква ва | ква стхулам ква ча ва сукшмам свамахімні стхітася ме ||}
\vperek{Де далеке, де близьке, де зовнішнє, де внутрішнє? Де грубе, де тонке — у мене, що стоїть у власній величі?}
\vcom{Построкове зіставлення: редакторська чернетка, не рецензовано.}
\vcom{Рядок а}
\vdeva{क्व दूरं क्व समीपं वा बाह्यं क्वाभ्यन्तरं क्व वा}
\vtranslit{ква дурам ква саміпам ва бахям квабгянтарам ква ва}
\vperek{Де далеке, де близьке, де зовнішнє, де внутрішнє?}
\vcom{dūraṃ / samīpaṃ позначають далеке й близьке; bāhyaṃ / abhyantaraṃ — зовнішнє й внутрішнє. Це триває низка питань kva, а не твердження про зникнення фізичних відстаней. У GRETIL можливе написання bahyaṃ, у SARIT bāhyaṃ: перевірити довготу голосної при фаховій колації.}
\vcom{Рядок б}
\vdeva{क्व स्थूलं क्व च वा सूक्ष्मं स्वमहिम्नि स्थितस्य मे}
\vtranslit{ква стхулам ква ча ва сукшмам свамахімні стхітася ме}
\vperek{Де грубе, де тонке — у мене, що стоїть у власній величі?}
\vcom{sthūlaṃ та sūkṣmaṃ зіставляють грубе або відчутне з тонким; обидва стоять у питальній конструкції. svamahimni sthitasya me — той самий рефрен, що в 19.2–19.5; переклад «у власній величі» не змінено.}
\vcom{Три пари, і кожна описує простір по-своєму: відстань (далеке / близьке), напрямок (зовнішнє / внутрішнє), щільність (грубе / тонке). Разом вони мали б вичерпати всі способи бути десь — і саме тому всі три відкидаються разом. Пара sthūla / sūkṣma особлива: це не лише «велике й мале», а й «тілесне й тонке», тобто остання межа між видимим і невидимим. РЕФРЕН ГЛАВИ. 19.2–19.7 закінчуються однаково: svamahimni sthitasya me («у мене, що стоїть у власній величі»). Кожен із цих шести віршів — низка питань kva («де»), і жодне не дістає відповіді. 19.1 і 19.8 стоять поза серією: перший описує, ЯК це сталося, останній — ЩО з цього випливає для мови.}
\vpar{Українська паралель (наша інтерпретація)}{Володимир Свідзінський, «Медобір» (1975, посмертно).}{Наша паралель: у Свідзінського далеке й близьке теж втрачають міру, коли зникає той, хто міряє. Наше порівняння — за образом.}
\vfig{ilustratsii/19-6-neti.pdf}{Далеке й близьке, грубе й тонке відпали}{Схема наша, не історична ілюстрація.}

\vsec{19.7}
\vdeva{क्व मृत्युर् जीवितं वा क्व लोकाः क्वास्य क्व लौकिकम् । क्व लयः क्व समाधिर् वा स्वमहिम्नि स्थितस्य मे ॥}
\vtranslit{ква мрітюр джівітам ва ква локах квася ква лаукікам | ква лаях ква самадгір ва свамахімні стхітася ме ||}
\vperek{Де смерть, де життя, де світи, де мирське? Де розчинення, де самадгі — у мене, що стоїть у власній величі?}
\vcom{Построкове зіставлення: редакторська чернетка, не рецензовано.}
\vcom{Рядок а}
\vdeva{क्व मृत्युर् जीवितं वा क्व लोकाः क्वास्य क्व लौकिकम्}
\vtranslit{ква мрітюр джівітам ва ква локах квася ква лаукікам}
\vperek{Де смерть, де життя, де світи, де мирське?}
\vcom{mṛtyuḥ — «смерть», jīvitaṃ — «життя», lokāḥ — «світи», laukikam — «мирське». Форма kvāsya у чинному санскритському рядку граматично неясна й не має однозначного українського відповідника в чинній картці; її **не** можна непомітно тлумачити як встановлене читання.}
\vcom{Рядок б}
\vdeva{क्व लयः क्व समाधिर् वा स्वमहिम्नि स्थितस्य मे}
\vtranslit{ква лаях ква самадгір ва свамахімні стхітася ме}
\vperek{Де розчинення, де самадгі — у мене, що стоїть у власній величі?}
\vcom{layaḥ означає розчинення, samādhiḥ — зосередження або самадгі; kva ставить обидві практичні цілі під питання. Рефрен svamahimni sthitasya me зберігається; вірш не є доказом безкорисності будь-якої медитації поза власним філософським контекстом.}
\vcom{Останній із шести переліків, і він забирає найвище: laya і samādhi — те, до чого глави 5 і 13 кликали, тут теж не мають місця. Це замикає книгу на собі: навіть мета попередніх глав не стає притулком. Форма kvāsya (kva + asya, «де цього») темна: asya — родового від займенника, але без опори в реченні; можливо, це зіпсоване читання. У перекладі цей член пропущено, і це позначено, а не замасковано. РЕФРЕН ГЛАВИ. 19.2–19.7 закінчуються однаково: svamahimni sthitasya me («у мене, що стоїть у власній величі»). Кожен із цих шести віршів — низка питань kva («де»), і жодне не дістає відповіді. 19.1 і 19.8 стоять поза серією: перший описує, ЯК це сталося, останній — ЩО з цього випливає для мови.}
\vpar{Українська паралель (наша інтерпретація)}{Леся Українка, «Кассандра» (1907).}{Наша паралель: у Лесі Українки життя й смерть теж не утворюють пари, вартої вибору, для того, хто вже стоїть поза нею. Наше порівняння — типологічне.}
\vfig{ilustratsii/19-7-neti.pdf}{Смерть, життя й самадгі відпали}{Схема наша, не історична ілюстрація.}

\vsec{19.8}
\vdeva{अलं त्रिवर्गकथया योगस्य कथयाप्य् अलम् । अलं विज्ञानकथया विश्रान्तस्य ममात्मनि ॥}
\vtranslit{алам тріварґакатхая йоґася катхаяпй алам | алам віджнянакатхая вішрантася маматмані ||}
\vperek{Досить розмов про трицінність, досить і розмов про йогу, досить розмов про знання — у мене, що спочив в Атмані.}
\vcom{Построкове зіставлення: редакторська чернетка, не рецензовано.}
\vcom{Рядок а}
\vdeva{अलं त्रिवर्गकथया योगस्य कथयाप्य् अलम्}
\vtranslit{алам тріварґакатхая йоґася катхаяпй алам}
\vperek{Досить розмов про трицінність, досить і розмов про йогу,}
\vcom{Повтор alam («досить») охоплює trivarga («три групи цілей» — дгарма, користь, бажання) і розмову про йогу; kathayā стоїть в орудному відмінку. Це підсумкова риторична формула, а не наказ заборонити українській спільноті обговорення знань.}
\vcom{Рядок б}
\vdeva{अलं विज्ञानकथया विश्रान्तस्य ममात्मनि}
\vtranslit{алам віджнянакатхая вішрантася маматмані}
\vperek{досить розмов про знання — у мене, що спочив в Атмані.}
\vcom{vijñānakathayā — «розмовою про знання». viśrāntasya mamātmani означає стан мовця, що спочив в Атмані; завершення відрізняється від рефрену svamahimni sthitasya me у віршах 19.2–19.7. Третє alam підкреслює завершальний характер вислову.}
\vcom{Вірш не має рефрену попередніх шести й замінює його іншим: не svamahimni sthitasya («що стоїть у власній величі»), а viśrāntasya mamātmani («що спочив в Атмані»). Різниця точна: стояння — це позиція, спочинок — її відсутність потреби. Тричі повторене alam («досить») знищує три теми поспіль: етику (трицінність), практику (йога) і теорію (знання). Це останній вірш книги перед главой 20, і він закриває розмову — буквально: kathā означає «розмова», і саме розмову тут оголошено зайвою.}
\vpar{Українська паралель (наша інтерпретація)}{Григорій Сковорода, «Розмова пʼяти подорожніх про істинне щастя в житті».}{Наша паралель: у Сковороди розмови про знання теж урешті виявляються зайвими перед тим, що вже є. Наше порівняння — типологічне.}
\vfig{ilustratsii/19-8-flame.pdf}{Полумʼя, що втихло: досить розмов}{Схема наша, не історична ілюстрація.}

